
\documentclass{article}
\usepackage[style=ieee]{biblatex}
\addbibresource{synthetic_cs_engineering_ieee_clean.bib}

\title{A Synthetic Manuscript for Reference-Extraction Benchmarking}
\author{Synthetic Author}
\date{2026}

\begin{document}
\maketitle

\begin{abstract}
This synthetic manuscript exists only to exercise the reference-extraction pipeline against a document with realistic structure and an exactly known bibliography, generated from the entries in entries.py.
\end{abstract}

\section{Introduction}
The question taken up here has been approached from several directions, and the paragraphs below trace the line of work this manuscript builds on rather than restating each contribution in full. A closely comparable argument is developed by \textcite{shannon1948mathematical}. This reading sits comfortably beside the surrounding literature \parencite{vaswani2017attention}. The argument is put most directly in a single passage \parencite[p.~14]{rcoreteam2024r}. The finding has since been replicated more than once \parencite{turing1950computing,codd1970relational}. This reading sits comfortably beside the surrounding literature \parencite{lamport1978time}. The argument is put most directly in a single passage \parencite[p.~14]{rivest1978method}. The finding has since been replicated more than once \parencite{lecun1998gradient,diffiehellman1976new}.

\section{Method}
The procedure follows established practice at every step, so what matters in this section is which earlier report each decision was taken from. Earlier evidence points in this direction as well \parencite{krizhevsky2012imagenet}. The relevant table is reproduced there \parencite[p.~87]{dijkstra1968goto}. Two independent reports agree on this point \parencite{knuth1968art,church1936unsolvable}. Earlier evidence points in this direction as well \parencite{vonneumann1945firstdraft}. The relevant table is reproduced there \parencite[p.~87]{bernerslee1994worldwideweb}.

\section{Results}
The pattern observed in these data is reported first and then placed against the findings it is meant to be read alongside. The finding has since been replicated more than once \parencite{cerfkahn1974protocol,metropolis1953equation}. This reading sits comfortably beside the surrounding literature \parencite{hopcrofttarjan1973efficient}. The argument is put most directly in a single passage \parencite[p.~14]{cormen2009introduction}. The finding has since been replicated more than once \parencite{hopper1952education,engelbart1968research}. This reading sits comfortably beside the surrounding literature \parencite{ritchiethompson1974unix}. The argument is put most directly in a single passage \parencite[p.~14]{backus1978can}.

\section{Discussion}
What remains is to say which of the earlier claims the present data support and which of them they leave open. The finding has since been replicated more than once \parencite{mccarthy1960recursive,minskypapert1969perceptrons}. This reading sits comfortably beside the surrounding literature \parencite{rumelhart1986learning}. The argument is put most directly in a single passage \parencite[p.~14]{hochreiterschmidhuber1997long}. The finding has since been replicated more than once \parencite{goodfellow2014generative,pagebrin1999pagerank}. This reading sits comfortably beside the surrounding literature \parencite{tanenbaum1981computer}.

\printbibliography

\end{document}
